\documentclass[aspectratio=169,11pt]{beamer}
\usepackage[utf8]{inputenc}
\usepackage[T1]{fontenc}
\usepackage[spanish,es-nodecimaldot]{babel}
\usepackage{graphicx}
\usepackage{amsmath}
\usepackage{amssymb}
\usepackage{booktabs}
\usepackage{xcolor}

% Definición de Colores USS
\definecolor{USSBlue}{RGB}{0,32,91}
\definecolor{USSGold}{RGB}{212,175,55}
\definecolor{LightBlue}{RGB}{235,240,250}
\definecolor{DarkGray}{RGB}{80,80,80}


% Configuración del Tema Beamer
\usetheme{Madrid}
\usecolortheme{default}

% Personalización de Colores en Madrid
\setbeamercolor{structure}{fg=USSBlue}
\setbeamercolor{palette primary}{bg=USSBlue,fg=white}
\setbeamercolor{palette secondary}{bg=USSBlue!85!black,fg=white}
\setbeamercolor{palette tertiary}{bg=USSGold,fg=USSBlue!30!black}
\setbeamercolor{titlelike}{parent=palette primary}
\setbeamercolor{block title}{bg=USSBlue,fg=white}
\setbeamercolor{block body}{bg=LightBlue,fg=black}

% Configuración de las viñetas
\setbeamertemplate{items}[circle]

% Título y datos de portada
\title[Brazo Robótico 3 GDL]{\textbf{Diseño, Modelamiento Cinemático y Análisis de Integridad de un Brazo Robótico Antropomórfico de 3 GDL}}
\subtitle{Hito Cami - Taller de Aptitudes Lógicas y Matemáticas}
\author[Moisés, Fernando, Juan Pablo]{Moisés, Fernando, Juan Pablo}
\institute[USS]{
    Facultad de Ingeniería y Tecnología\\
    Carrera de Ingeniería Civil\\
    \textbf{Universidad San Sebastián}
}
\date[Junio 2026]{Junio, 2026}

% Insertar logo en la portada
\titlegraphic{
\includegraphics[height=1.2cm]{logo_uss.png}}

\begin{document}

% Slide 1: Portada
\begin{frame}
    \titlepage
\end{frame}

% Slide 2: Contexto y Analogía del Brazo Humano
% [VALIDACIÓN LÓGICA]: Satisface el criterio de "Comunicación Matemática" (CM) de la presentación al introducir claramente el contexto y objetivos del manipulador.
\begin{frame}{Contexto y Analogía del Brazo Humano}
    \begin{columns}
        \begin{column}{0.55\textwidth}
            El brazo robótico antropomórfico de 3 GDL imita la estructura de la extremidad superior humana:
            \vspace{0.3cm}
            \begin{itemize}
                \item \textbf{Cintura (Yaw):} Rotación de la base alrededor del eje vertical $z$. Ángulo $\theta_1$.
                \item \textbf{Hombro (Pitch):} Elevación principal sobre el eje horizontal. Ángulo $\theta_2$.
                \item \textbf{Codo (Pitch relativo):} Flexo-extensión distal respecto al primer eslabón. Ángulo $\theta_3$.
            \end{itemize}
            \vspace{0.3cm}
            \textbf{Objetivo:} Mapear analíticamente el espacio articular al espacio operacional tridimensional y evaluar la resistencia estructural ante rotación crítica.
        \end{column}
        \begin{column}{0.4\textwidth}
            \begin{block}{Grados de Libertad (GDL)}
                \centering
                \begin{itemize}
                    \item $\theta_1$: Yaw base
                    \item $\theta_2$: Pitch hombro
                    \item $\theta_3$: Pitch codo
                \end{itemize}
            \end{block}
            \centering
            
\includegraphics[width=0.6\textwidth]{logo_uss.png}
        \end{column}
    \end{columns}
\end{frame}

% Slide 3: Parámetros del Sistema y Supuestos
% [VALIDACIÓN LÓGICA]: Satisface el criterio de "Variables y Parámetros" (MM) identificando claramente las entradas geométricas, de masa y propiedades de los materiales con sus respectivas unidades físicas.
\begin{frame}{Variables y Parámetros del Sistema}
    \begin{columns}
        \begin{column}{0.5\textwidth}
            \begin{block}{Parámetros Geométricos}
                \begin{itemize}
                    \item Eslabones: $L_1 = 10\text{ cm}$, $L_2 = 10\text{ cm}$.
                    \item Alcance máximo radial: $R_{\text{max}} = 20\text{ cm} = 0.2\text{ m}$.
                \end{itemize}
            \end{block}
            \begin{block}{Parámetros Inerciales}
                \begin{itemize}
                    \item Masa eslabones: $m_1 = m_2 = 0.5\text{ kg}$.
                    \item Masa carga de pago: $M = 0.5\text{ kg}$.
                    \item Densidad lineal de eslabones: $\rho = 5.0\text{ kg/m}$.
                \end{itemize}
            \end{block}
        \end{column}
        \begin{column}{0.5\textwidth}
            \begin{block}{Materiales Seleccionados}
                \begin{itemize}
                    \item \textbf{Eslabones (Aluminio 6061-T6):}
                        \begin{itemize}
                            \item Fluencia: $\sigma_y = 276\text{ MPa}$
                            \item Expansión: $\alpha_{\text{Al}} = 23 \times 10^{-6}\text{ /}^\circ\text{C}$
                        \end{itemize}
                    \item \textbf{Pasador (Acero Inoxidable 316):}
                        \begin{itemize}
                            \item Cortante admisible: $\tau_{\text{adm}} = 150\text{ MPa}$
                            \item Expansión: $\alpha_{\text{Acero}} = 16 \times 10^{-6}\text{ /}^\circ\text{C}$
                        \end{itemize}
                \end{itemize}
            \end{block}
        \end{column}
    \end{columns}
\end{frame}

% Slide 4: Ecuaciones Cinemáticas (Analogía Sombra Radial)
% [VALIDACIÓN LÓGICA]: Satisface el criterio de "Uso de Representaciones" (CM) y "Limitaciones del Modelo" (RA) al derivar las ecuaciones del espacio operacional 3D y advertir sobre la singularidad geométrica de base.
\begin{frame}{Ecuaciones Cinemáticas: Analogía Sombra Radial}
    Proyectamos el brazo sobre el plano horizontal para hallar la sombra radial $r$:
    \begin{equation}
        r = L_1\cos(\theta_2) + L_2\cos(\theta_2 + \theta_3)
    \end{equation}
    
    A partir del radio proyectado $r$, las coordenadas cartesianas $(x, y, z)$ son:
    \begin{align}
        x &= r\cos(\theta_1) = \left[ L_1\cos(\theta_2) + L_2\cos(\theta_2 + \theta_3) \right]\cos(\theta_1) \\
        y &= r\sin(\theta_1) = \left[ L_1\cos(\theta_2) + L_2\cos(\theta_2 + \theta_3) \right]\sin(\theta_1) \\
        z &= L_1\sin(\theta_2) + L_2\sin(\theta_2 + \theta_3)
    \end{align}
    
    \begin{block}{Singularidad Geométrica de Base}
        Cuando $r = 0$ (brazo en posición vertical), $\theta_1$ queda indeterminado, provocando una singularidad del manipulador en coordenadas cartesianas.
    \end{block}
\end{frame}

% Slide 5: Espacio de Trabajo Tridimensional
\begin{frame}{Espacio de Trabajo Tridimensional (Workspace)}
    \begin{columns}
        \begin{column}{0.5\textwidth}
            La envolvente tridimensional del brazo está restringida por la longitud física y los rangos angulares:
            \vspace{0.2cm}
            \begin{itemize}
                \item Envolvente máxima teórica:
                $$x^2 + y^2 + z^2 \le 400\text{ cm}^2 \quad (\text{Radio } 20\text{ cm})$$
                \item Límites angulares del diseño:
                \begin{itemize}
                    \item $\theta_1 \in [-180^\circ, 180^\circ]$
                    \item $\theta_2 \in [-90^\circ, 90^\circ]$
                    \item $\theta_3 \in [-150^\circ, 150^\circ]$
                \end{itemize}
            \end{itemize}
        \end{column}
        \begin{column}{0.5\textwidth}
            \centering
            \includegraphics[width=\textwidth]{workspace_3d.png}
        \end{column}
    \end{columns}
\end{frame}

% Slide 6: Escenario Crítico de Rotación (140 km/h)
% [VALIDACIÓN LÓGICA]: Satisface el criterio de "Construcción del Modelo" (MM) al formular la integral dinámica de tracción centrípeta para modelar esfuerzos bajo una velocidad crítica de 140 km/h.
\begin{frame}{Escenario Crítico de Rotación (140 km/h)}
    Evaluamos el brazo totalmente extendido ($R_{\text{max}} = 0.2\text{ m}$) a $140\text{ km/h}$ ($38.89\text{ m/s}$):
    \vspace{0.2cm}
    \begin{columns}
        \begin{column}{0.5\textwidth}
            \begin{block}{Parámetros Cinemáticos}
                \begin{itemize}
                    \item Velocidad angular: $\omega = \frac{v}{R_{\text{max}}} = 194.44\text{ rad/s}$
                    \item Aceleración centrípeta en extremo:
                    $$a_c = \omega^2 R_{\text{max}} = 7561.73\text{ m/s}^2 \approx \mathbf{771g}$$
                \end{itemize}
            \end{block}
        \end{column}
        \begin{column}{0.5\textwidth}
            \begin{block}{Cálculo de Tensiones Centrífugas}
                La tensión a lo largo del radio $x$ es:
                $$T(x) = \int_x^{L_1+L_2} \rho \omega^2 \xi d\xi + M \omega^2 (L_1+L_2)$$
                \begin{itemize}
                    \item \textbf{Base del brazo ($x=0$):}
                    $$T_{\text{base}} = \mathbf{7561.73\text{ N}}$$
                    \item \textbf{Articulación del codo ($x=0.1\text{ m}$):}
                    $$T_{\text{codo}} = \mathbf{6616.51\text{ N}}$$
                \end{itemize}
            \end{block}
        \end{column}
    \end{columns}
\end{frame}

% Slide 7: Verificación de Esfuerzos y Seguridad
% [VALIDACIÓN LÓGICA]: Satisface el criterio de "Defensa de Decisiones" al calcular cuantitativamente los factores de seguridad del aluminio de los eslabones y el acero del pasador, identificando el punto de falla del diseño original.
\begin{frame}{Verificación de Esfuerzos y Seguridad}
    \begin{columns}
        \begin{column}{0.5\textwidth}
            \begin{block}{Eslabón de Aluminio (Base)}
                \begin{itemize}
                    \item Diámetro: $10\text{ mm}$ (Área $78.54\text{ mm}^2$)
                    \item Tensión: $T_{\text{base}} = 7561.73\text{ N}$
                    \item Esfuerzo normal:
                    $$\sigma_{\text{link}} = \frac{T_{\text{base}}}{A_{\text{link}}} = \mathbf{96.28\text{ MPa}}$$
                    \item Factor de Seguridad (Fluencia $276\text{ MPa}$):
                    $$FS_{\text{link}} = \frac{276}{96.28} = \mathbf{2.87} \quad (\textbf{RESISTE})$$
                \end{itemize}
            \end{block}
        \end{column}
        \begin{column}{0.5\textwidth}
            \begin{block}{Pasador Original de Acero (Codo)}
                \begin{itemize}
                    \item Diámetro: $3\text{ mm}$ (Área $7.07\text{ mm}^2$, cizalle simple)
                    \item Fuerza de corte: $T_{\text{codo}} = 6616.51\text{ N}$
                    \item Esfuerzo cortante:
                    $$\tau_{\text{pin,3mm}} = \frac{T_{\text{codo}}}{A_{\text{pin,3mm}}} = \mathbf{936.0\text{ MPa}}$$
                    \item Factor de Seguridad (Admisible $150\text{ MPa}$):
                    $$FS_{\text{pin,3mm}} = \frac{150}{936.0} = \mathbf{0.16} \quad (\textbf{FALLA CRÍTICA})$$
                \end{itemize}
            \end{block}
        \end{column}
    \end{columns}
\end{frame}

% Slide 8: Análisis de Robustez Dinámica y Térmica
% [VALIDACIÓN LÓGICA]: Satisface el criterio de "Prueba de Robustez" (RA) y "Limitaciones del Modelo" (RA) al tabular la sensibilidad de esfuerzos frente a la velocidad y detallar el comportamiento del juego del codo por dilatación diferencial.
\begin{frame}{Análisis de Robustez Dinámica y Térmica}
    \begin{columns}
        \begin{column}{0.55\textwidth}
            \begin{table}[H]
                \centering
                \caption{Esfuerzos vs. Velocidad Tangencial.}
                \resizebox{\textwidth}{!}{
                \begin{tabular}{ccccc}
                    \toprule
                    \textbf{Velocidad} & \textbf{Acel. Centr.} & \textbf{Tensión Codo} & \textbf{$\tau_{\text{pin,3mm}}$} & \textbf{Estado Pin} \\
                    \midrule
                    40 km/h & $62.9g$ & $540.1\text{ N}$ & $76.4\text{ MPa}$ & Resiste \\
                    \textbf{56 km/h} & $123.3g$ & $1060.3\text{ N}$ & $150.0\text{ MPa}$ & \textbf{Límite} \\
                    80 km/h & $251.8g$ & $2160.5\text{ N}$ & $305.6\text{ MPa}$ & Falla \\
                    \textbf{140 km/h} & $771.6g$ & $6616.5\text{ N}$ & $936.0\text{ MPa}$ & \textbf{Falla Crít.} \\
                    200 km/h & $1573.6g$ & $13503.1\text{ N}$ & $1910.2\text{ MPa}$ & Falla Sev. \\
                    \bottomrule
                \end{tabular}
                }
            \end{table}
        \end{column}
        \begin{column}{0.45\textwidth}
            \begin{block}{Dilatación Térmica ($\Delta T = 100^\circ$C)}
                \begin{itemize}
                    \item \textbf{Eslabones:} Alargamiento de $0.23\text{ mm}$ cada uno (alarga el alcance radial en $0.46\text{ mm}$).
                    \item \textbf{Juego del codo:} El orificio de Al se expande más que el pasador de Acero ($\Delta d_{\text{orificio}} = 18.4\,\mu\text{m}$, $\Delta d_{\text{pin}} = 12.8\,\mu\text{m}$).
                    \item \textbf{Resultado:} Juego neto aumenta en $\mathbf{5.6\,\mu\text{m}}$. Previene el engripamiento.
                \end{itemize}
            \end{block}
        \end{column}
    \end{columns}
\end{frame}

% Slide 9: Propuesta de Rediseño de la Articulación
% [VALIDACIÓN LÓGICA]: Satisface el criterio de "Construcción del Modelo" (MM) al demostrar cuantitativamente la efectividad de la alternativa de rediseño de pasador de 8 mm en cortante doble.
\begin{frame}{Propuesta de Rediseño de la Articulación del Codo}
    \begin{columns}
        \begin{column}{0.5\textwidth}
            \textbf{Propuesta:} Aumentar el diámetro del pasador a $8\text{ mm}$ e implementar una junta en cortante doble.
            \vspace{0.2cm}
            \begin{itemize}
                \item \textbf{Cortante Simple con $d = 8\text{ mm}$ (Área $50.27\text{ mm}^2$):}
                $$\tau_{\text{pin,8mm}} = \frac{T_{\text{codo}}}{A_{\text{pin,8mm}}} = \mathbf{131.63\text{ MPa}} < 150\text{ MPa}$$
                $$FS_{\text{pin}} = 1.14 \quad (\textbf{RESISTE})$$
                \item \textbf{Cortante Doble con $d = 8\text{ mm}$ (Área $2 \times 50.27\text{ mm}^2$):}
                $$\tau_{\text{pin,doble}} = \frac{T_{\text{codo}}}{2 A_{\text{pin}}} = \mathbf{65.81\text{ MPa}} \ll 150\text{ MPa}$$
                $$FS_{\text{pin,doble}} = \mathbf{2.28} \quad (\textbf{ALTAMENTE SEGURO})$$
            \end{itemize}
        \end{column}
        \begin{column}{0.5\textwidth}
            \centering
            \includegraphics[width=0.9\textwidth]{stress_analysis.png}
        \end{column}
    \end{columns}
\end{frame}

% Slide 10: Conclusiones y Trabajo en Equipo
% [VALIDACIÓN LÓGICA]: Satisface el criterio de "Conclusión y Reflexión" (RP+RA+CM) al sintetizar los hallazgos mecánicos clave y explicitar los roles y contribuciones individuales del equipo.
\begin{frame}{Conclusiones y Trabajo en Equipo}
    \begin{columns}
        \begin{column}{0.5\textwidth}
            \begin{block}{Conclusiones Clave}
                \begin{itemize}
                    \item Los eslabones de Al 6061-T6 resisten a 140 km/h ($FS=2.87$).
                    \item El pasador original de 3 mm falla inmediatamente a 140 km/h ($FS=0.16$).
                    \item Límite de velocidad del diseño original: \textbf{56 km/h}.
                    \item El rediseño a pasador de 8 mm en cortante doble es óptimo ($FS=2.28$).
                    \item La variación térmica de $100^\circ$C genera $5.6\,\mu\text{m}$ de juego libre adicional sin agarrotamiento.
                \end{itemize}
            \end{block}
        \end{column}
        \begin{column}{0.5\textwidth}
            \begin{block}{Trabajo en Equipo (Roles)}
                \begin{itemize}
                    \item \textbf{Moisés:} Cinemática directa 3D y análisis de dinámica centrífuga.
                    \item \textbf{Fernando:} Verificación estructural de esfuerzos y pasador.
                    \item \textbf{Juan Pablo:} Modelado termoelástico y validaciones numéricas.
                \end{itemize}
            \end{block}
            \begin{block}{Herramientas de Apoyo}
                Ecosistema de co-diseño con \textbf{Obsidian} para control de notas y \textbf{OpenCode} para la automatización y simulación en Python.
            \end{block}
        \end{column}
    \end{columns}
\end{frame}

% Slide 11: Preguntas y Respuestas (Q&A)
\begin{frame}
    \centering
    \Huge \textcolor{USSBlue}{\textbf{Preguntas y Respuestas (Q\&A)}}\\[1cm]
    \large \textcolor{USSBlue}{¡Muchas gracias por su atención!}\\[0.5cm]
    \small \textcolor{DarkGray}{Taller de Aptitudes Lógicas y Matemáticas -- Hito Cami}\\[0.5cm]
    
\includegraphics[height=1.2cm]{logo_uss.png}
\end{frame}

\end{document}
